\documentclass[handout]{beamer}

\usepackage[utf8]{inputenc}
\usepackage[T1]{fontenc}
\usepackage[brazil]{babel}
\usepackage{listings}

\usetheme{Madrid}

\title{\LaTeX\ para gente grande}
\subtitle{Dia 1: Ecossistema, Compilação e Fontes}
\author[Andrês Oliveira]{Andrês Oliveira}
\date{Minicurso LEM/UNICAMP}

\definecolor{codegreen}{rgb}{0,0.6,0}
\definecolor{codegray}{rgb}{0.5,0.5,0.5}
\definecolor{codepurple}{rgb}{0.58,0,0.82}
\definecolor{backcolour}{rgb}{0.95,0.95,0.92}

\lstset{
    backgroundcolor=\color{backcolour},   
    commentstyle=\color{codegreen},
    keywordstyle=\color{magenta},
    numberstyle=\tiny\color{codegray},
    stringstyle=\color{codepurple},
    basicstyle=\ttfamily\scriptsize,
    breaklines=true,                 
    showstringspaces=false,
    tabsize=2
}

\begin{document}

% BLOCO 1 - ABERTURA
\begin{frame}
    \maketitle
\end{frame}

\begin{frame}{Objetivos da Tarde}
    \begin{itemize}
        \item Dominar a terminologia profissional do ecossistema TeX.
        \item Compreender o fluxo de dados entre arquivos auxiliares.
        \item Configurar fontes no \texttt{pdfLaTeX} para o português.
        \item Diagnosticar erros de compilação sistematicamente.
        \item Conhecer ferramentas de automação (\texttt{latexmk}, \texttt{make}).
    \end{itemize}
\end{frame}

\begin{frame}{O que você saberá ao final do dia}
    \begin{itemize}
        \item Por que o sumário às vezes some e como recuperá-lo.
        \item Como ler o log sem medo de mensagens criptografadas.
        \item A diferença técnica entre codificações T1 e OT1.
        \item Como consultar manuais oficiais via terminal.
    \end{itemize}
\end{frame}

% BLOCO 2 - ECOSSISTEMA
\begin{frame}{LaTeX não é um programa único}
    \begin{block}{O Ecossistema TeX}
        LaTeX é um conjunto de macros rodando sobre um motor tipográfico, integrado em um condomínio de ferramentas.
    \end{block}
    \vfill
    \textbf{Principais componentes:}
    \begin{enumerate}
        \item Distribuição (O repositório).
        \item Compilador/Motor (O cérebro).
        \item Editor (A interface).
        \item Classe e Pacote (As regras e funções).
    \end{enumerate}
\end{frame}

\begin{frame}{Distribuições: O Repositório de Ferramentas}
    \begin{itemize}
        \item \textbf{TeX Live}: O padrão ouro multiplataforma (Linux, Mac, Windows).
        \item \textbf{MiKTeX}: Popular no Windows, instala pacotes sob demanda.
        \item \textbf{MacTeX}: A versão otimizada do TeX Live para macOS.
    \end{itemize}
    \vfill
    \begin{block}{Observação técnica}
        Ambientes online como o \textbf{Overleaf} utilizam distribuições TeX Live nos bastidores.
    \end{block}
\end{frame}

\begin{frame}{Editores: Onde a mágica (não) acontece}
    O editor \textbf{não é} o compilador. Ele apenas chama o programa que faz o trabalho.
    \begin{itemize}
        \item \textbf{TeXstudio}: IDE completa com visualizadores de log e PDF integrados.
        \item \textbf{VS Code + LaTeX Workshop}: Fluxo moderno e altamente customizável.
        \item \textbf{Overleaf}: Editor na nuvem (ideal para colaboração em tempo real).
        \item \textbf{Emacs + AUCTeX}: O auge da automação para usuários avançados.
    \end{itemize}
\end{frame}

\begin{frame}{Compilador (Motor): Quem processa}
    No contexto técnico, o compilador é frequentemente chamado de \textbf{Engine} ou motor.
    \begin{itemize}
        \item \textbf{\texttt{pdflatex}}: Estável, rápido, foco em fontes tradicionais.
        \item \textbf{\texttt{xelatex}}: Unicode nativo, usa fontes do sistema (.ttf, .otf).
        \item \textbf{\texttt{lualatex}}: Extensível via scripts Lua, suporte a fontes modernas.
    \end{itemize}
\end{frame}

\begin{frame}{Classe e Pacote}
    \begin{itemize}
        \item \textbf{Classe (\texttt{.cls})}: O esqueleto.
              \begin{itemize}
                  \item \texttt{article}, \texttt{book}, \texttt{beamer}, \texttt{abntex2}.
              \end{itemize}
        \item \textbf{Pacote (\texttt{.sty})}: Os superpoderes.
              \begin{itemize}
                  \item \texttt{amsmath} (matemática), \texttt{graphicx} (imagens), \texttt{tikz} (desenho).
              \end{itemize}
    \end{itemize}
\end{frame}

% BLOCO 3 - FLUXO DE COMPILAÇÃO
\begin{frame}{Anatomia da Compilação}
    \begin{center}
        \texttt{arquivo.tex} $\to$ \textbf{Compilador} $\to$ \texttt{arquivo.pdf}
    \end{center}
    \vfill
    \textbf{Arquivos Auxiliares (A Memória):}
    \begin{itemize}
        \item \textbf{\texttt{.aux}}: Onde o TeX anota labels e referências.
        \item \textbf{\texttt{.log}}: Onde o TeX anota tudo o que aconteceu (erros e avisos).
        \item \textbf{\texttt{.toc}}: Onde reside a estrutura do sumário.
        \item \textbf{\texttt{.bbl}}: A bibliografia formatada pelo BibTeX.
    \end{itemize}
\end{frame}

\begin{frame}{O Ciclo das Referências}
    Por que compilar mais de uma vez?
    \begin{enumerate}
        \item \textbf{1ª Passagem}: ``Vi um label chamado 'eq1' na página 5. Vou anotar no .aux''.
        \item \textbf{2ª Passagem}: ``Opa, alguém citou o 'eq1'. Vou ler no .aux onde ele estava''.
    \end{enumerate}
    \vfill
    \begin{alertblock}{Aviso clássico}
        \texttt{LaTeX Warning: Label(s) may have changed. Rerun to get cross-references right.}
    \end{alertblock}
\end{frame}

% BLOCO 4 - LABORATÓRIO 1
\begin{frame}[fragile]{Laboratório 1: O Vazio e o Erro}
    \textbf{Objetivo:} Perder o medo das mensagens de erro iniciais.

    \textbf{Tarefa:}
    \begin{enumerate}
        \item Criar um arquivo \texttt{vazio.tex} sem nenhum conteúdo.
        \item Tentar compilar: \texttt{pdflatex vazio.tex}.
        \item Observar a mensagem final.
    \end{enumerate}

    \textbf{Observe:} O compilador para em \texttt{*} esperando comandos ou reclama de \texttt{No pages of output}. Use \texttt{X} ou \texttt{Ctrl+C} para sair.
\end{frame}

\begin{frame}[fragile]{Laboratório 2: Primeiro Documento e Log}
    \textbf{Objetivo:} Identificar a anatomia básica de um projeto.

    \textbf{Tarefa:}
    \begin{enumerate}
        \item Criar o arquivo \texttt{ola.tex}:
\begin{lstlisting}[language={[latex]TeX}]
\documentclass{article}
\begin{document}
Ola, mundo!
\end{document}
\end{lstlisting}
        \item Compilar e rodar \texttt{ls} no terminal.
        \item Abrir o arquivo \texttt{ola.log} e procurar pela palavra ``pdfTeX''.
    \end{enumerate}

    \textbf{Resultado:} Você verá arquivos .aux e .log. No log, o compilador lista a versão do motor e pacotes usados.
\end{frame}

% BLOCO 5 - MOTORES E FONTES
\begin{frame}{pdflatex vs Motores Modernos}
    \begin{itemize}
        \item \textbf{pdflatex}: Padrão de estabilidade. Usa o modelo tradicional do TeX.
        \item \textbf{fontspec}: Pacote para Xe/LuaTeX que carrega fontes do sistema.
    \end{itemize}
    \vfill
    \begin{block}{Nossa escolha: pdflatex}
        Garante reprodutibilidade sem depender de fontes instaladas no sistema operacional do colega.
    \end{block}
\end{frame}

\begin{frame}[fragile]{Configuração de Fontes em Português}
    No \texttt{pdfLaTeX}, o segredo está no tripé:
    \begin{itemize}
        \item \textbf{UTF-8}: Codificação de entrada.
        \item \textbf{T1}: Codificação de saída (glifos acentuados únicos).
        \item \textbf{Pacote}: A família tipográfica.
    \end{itemize}
    \vfill
\begin{lstlisting}[language={[latex]TeX}]
\usepackage[utf8]{inputenc}
\usepackage[T1]{fontenc}
\usepackage{ebgaramond}
\end{lstlisting}
\end{frame}

\begin{frame}{Por que o T1 é inegociável?}
    \begin{itemize}
        \item \textbf{Sem T1}: ``ç'' vira ``c'' + ``acento'', quebra a hifenação e a busca no PDF.
        \item \textbf{Com T1}: O PDF contém o glifo real de ``ç''.
    \end{itemize}
    \vfill
    \textit{Teste no PDF final: tente copiar uma palavra com acento. Se sair com lixo, o T1 está faltando.}
\end{frame}

\begin{frame}{Microtipografia: O Refinamento Final}
    O pacote \texttt{microtype} faz ajustes microscópicos automáticos.
    \begin{itemize}
        \item \textbf{Protrusão}: Alinha pontuações na margem direita.
        \item \textbf{Expansão}: Ajusta largura das letras para evitar hifenação.
    \end{itemize}
    \vfill
    Torna blocos de texto longo muito mais agradáveis de ler.
\end{frame}

\begin{frame}[fragile]{fontspec (Apenas para conhecimento)}
    \textbf{Como seria em XeLaTeX/LuaLaTeX:}
\begin{lstlisting}[language={[latex]TeX}]
\usepackage{fontspec}
\setmainfont{Arial} % Usa fonte do sistema
\end{lstlisting}
    \vfill
    \begin{alertblock}{Importante}
        Este código \textbf{não funciona} no nosso projeto porque usamos \texttt{pdflatex}.
    \end{alertblock}
\end{frame}

% BLOCO 6 - AUTOMAÇÃO
\begin{frame}{O cansaço da compilação manual}
    Para um projeto com bibliografia:
    \begin{enumerate}
        \item \texttt{pdflatex doc}
        \item \texttt{bibtex doc}
        \item \texttt{pdflatex doc}
        \item \texttt{pdflatex doc}
    \end{enumerate}
    \vfill
    \textbf{Fadiga técnica}: Esquecer um passo causa referências quebradas.
\end{frame}

\begin{frame}{Ferramentas de Automação}
    \begin{itemize}
        \item \textbf{\texttt{latexmk}}: Script inteligente que monitora dependências e re-compila o necessário.
        \item \textbf{Makefile}: Padroniza o build. No curso, basta digitar \texttt{make}.
        \item \textbf{Editor/Overleaf}: Já fazem isso via botões, mas escondem o processo.
    \end{itemize}
\end{frame}

% BLOCO 7 - ERROS E DIAGNÓSTICO
\begin{frame}{Estratégia de Diagnóstico}
    \begin{alertblock}{Regra de Ouro}
        Sempre ignore os últimos erros e leia o \textbf{primeiro} erro relevante.
    \end{alertblock}
    \vfill
    \textbf{Como ler o erro:}
    \begin{itemize}
        \item Localize o ponto de exclamação (\texttt{!}).
        \item Verifique o número da linha.
        \item Olhe a linha acima da indicação do erro.
    \end{itemize}
\end{frame}

\begin{frame}{Undefined control sequence}
    \textbf{Causa:} O TeX viu um comando que ele não conhece.
    \begin{itemize}
        \item Erro de digitação (\texttt{\textbackslash seciton}).
        \item Esqueceu de carregar um pacote (\texttt{\textbackslash includegraphics} sem \texttt{graphicx}).
    \end{itemize}
\end{frame}

\begin{frame}{Missing \$ inserted}
    \textbf{Causa:} Uso de caracteres reservados fora do ambiente matemático.
    \vfill
    Exemplos perigosos:
    \begin{itemize}
        \item \texttt{\_} (underline);
        \item \texttt{\textasciicircum} (circunflexo);
        \item Símbolos matemáticos digitados diretamente no texto.
    \end{itemize}
\end{frame}

\begin{frame}{Overfull \textbackslash hbox (A caixa vazada)}
    \textbf{Sintoma:} Texto saindo da margem no PDF.
    \vfill
    \textbf{Causa:} O algoritmo do TeX preferiu estourar a margem do que criar espaços muito grandes entre as palavras.
    \vfill
    \textbf{Solução:} Verifique a hifenação (\texttt{babel} + \texttt{fontenc} T1) ou reformule a frase.
\end{frame}

% BLOCO 8 - PRÁTICA GUIADA
\begin{frame}[fragile]{Laboratório 3: Provocando o Caos}
    \textbf{Objetivo:} Localizar e corrigir um erro de comando.

    \textbf{Tarefa:}
    \begin{enumerate}
        \item No arquivo \texttt{ola.tex}, digite: \texttt{\textbackslash errinho\{Texto\}}.
        \item Tente compilar.
        \item Procure no terminal o erro \texttt{! Undefined control sequence}.
    \end{enumerate}

    \textbf{Verificação:} Apague o erro, recompile e veja se o arquivo .aux foi atualizado.
\end{frame}

\begin{frame}[fragile]{Laboratório 4: Referências Cruzadas}
    \textbf{Objetivo:} Observar a necessidade de múltiplas passagens.

    \textbf{Tarefa:}
    \begin{enumerate}
        \item Adicione um label: \texttt{\textbackslash section\{Intro\}\textbackslash label\{sec:intro\}}.
        \item No fim, cite: \texttt{Veja a sec \textbackslash ref\{sec:intro\}}.
        \item Compile \textbf{apenas uma vez} e veja o PDF.
    \end{enumerate}

    \textbf{Observe:} O PDF mostrará ``??''. Compile novamente e veja a mágica.
\end{frame}

\begin{frame}{Laboratório 5: Documentação local}
    \textbf{Objetivo:} Ganhar independência na resolução de problemas.

    \textbf{Tarefa:}
    \begin{enumerate}
        \item No terminal, digite: \texttt{texdoc geometry}.
        \item Tente procurar no PDF aberto como mudar as margens.
    \end{enumerate}

    \textbf{Resultado:} Você deve ter acesso ao manual oficial, que é mais preciso que tutoriais de blog.
\end{frame}

% BLOCO 9 - FECHAMENTO
\begin{frame}{Recapitulação}
    \begin{itemize}
        \item LaTeX é um ecossistema, não um editor.
        \item O compilador conversa com você via arquivo \texttt{.log}.
        \item Codificações corretas evitam 90\% dos problemas de português.
        \item \texttt{latexmk} e \texttt{make} são seus aliados.
    \end{itemize}
\end{frame}

\begin{frame}{Próximo Dia}
    \begin{itemize}
        \item Estruturação de documentos técnicos grandes.
        \item Matemática avançada com \texttt{amsmath} e \texttt{mathtools}.
        \item Cruzamento de dados e bibliografia profissional.
    \end{itemize}
\end{frame}

\end{document}
